%------------------------------------------------------------------------------%
\begin{exercises}

%------------------------------------------------------------------------------%
\begin{exercise}
Calcule as integrais usando a regra da substituição.
\begin{mathlist}
  \item \int x \sin(x^2) dx
  \item \int (x+1)\sqrt{2x+x^2}dx
  \item \int \frac{e^u}{(1+e^u)^2} du
  \item \int \cos^4(x) \sin(x) dx
  \item \int \frac{(\ln x)^2}{x} dx
  \item \int 5^t \sin(5^t)dt
  \item \int e^{\tan(x)} \sec^2(x)dx
  \item \int \frac{dx}{\sqrt{1-x^2} \arcsin(x)}
  \item \int \sqrt{cotg(x)} \csc^2(x) dx
  \item \int \sin(t) \sec^2(\cos(t))dt
  \item \int \frac{x}{1+x^4} dx
  \item \int x^3\sqrt{x^2+1} dx
  \item \int \frac{\arctan(x)}{1+x^2} dx
  \item \int \frac{dx}{ax+b}$, $a \neq 0
  \item \int_{0}^{1} \cos(\frac{\pi t}{2}) dt
  \item \int_0^2 \frac{e^{\frac{1}{x}}}{x^2} dx
  \item \int_{0}^{a} x(\sqrt{x^2+a^2})dx$, $a>0
  \item \int_0^1 \frac{dx}{(1+\sqrt{x})^4} \vert dx
  \item \int_0^4 \frac{x}{\sqrt{1+2x}} dx
  \item \int_e^{e^4} \frac{dx}{x \sqrt{\ln(x)}}
  \item \int_{\frac{1}{6}}^{\frac{1}{2}} \csc(\pi t) cotg(\pi t)dt
  \item \int_0^1 (3t-1)^{50} dt
  \item \int_0^1 \sqrt[3]{1+7x} dx
  \item \int_0^{13} \frac{dx}{\sqrt[3]{(1+2x)^2}} dt
  \item \int\frac{(x+1)^{2}\arctan(3x)+9x^{3}+x}{(9x^{2}+1)(x+1)^{2}} d x
  \item \int \sec^2(\theta)\tan^3(\theta) d\theta
\end{mathlist}
\begin{answer}
  \begin{mathlist}
  \item -\frac{1}{2} \cos(x^2)+C
  \item \frac{1}{3}(2x+x^2)^{\frac{3}{2}}+C
  \item \frac{1}{1-e^u}+C
  \item -\frac{1}{5} \cos^5(x)+C
  \item \frac{1}{2}(\ln(x))^2+C
  \item -\frac{1}{\ln(x)}\cos(5^t)+C
  \item e^{\tan(x)}+C
  \item \ln(\vert \arcsin(x) \vert)+ C
  \item -\frac{2}{3}(cotg(x))^{\frac{3}{2}}+C
  \item -\tan(\cos(x))+C
  \item \frac{1}{5}(x^2+1)^{\frac{5}{2}}- \frac{1}{3}(x^2+1)^{\frac{3}{2}}+C
  \item \frac{\arctan^2(x)}{2}+C
  \item \frac{1}{a}\ln(\vert ax+b \vert)
  \item \frac{2}{\pi}
  \item e-\sqrt{e}
  \item \frac{1}{3}(2\sqrt{2}-1)a^3
  \item -\frac{1}{6}
  \item \frac{10}{3}
  \item 2
  \item \frac{1}{\pi}
  \item \frac{1}{153}(2^{51}+1)
  \item \frac{45}{28}
  \item 3
  \item \frac{\arctan(3x)}{2}+ \ln(\vert x+1 \vert)+ \frac{1}{x+1}+K
  \item \frac{1}{4} \tan^4(\theta)+C
  \end{mathlist}
\end{answer}
\end{exercise}

%------------------------------------------------------------------------------%
\begin{exercise}
Calcule a área sob a curva dada no intervalo especificado.
\begin{mathlist}
  \item y= \arctan(x),                \quad 0 \leq x \leq 1
  \item y= \cos(x) \sin(\sin(x)),     \quad 0 \leq x \leq  \frac{\pi}{2}
  \item y= \frac{\sqrt{u}-2u^2}{u}du, \quad 1 \leq u \leq 9
\end{mathlist}
\begin{answer}
\begin{mathlist}
  \item \frac{1}{2}
  \item 1-\cos(1)
  \item \;
\end{mathlist}
\end{answer}
\end{exercise}

%------------------------------------------------------------------------------%
\begin{exercise}
Primeiro faça uma substituição e então use integração
por partes para calcular a integral
\begin{mathlist}
  \item \int \cos(\sqrt{x}) dx
  \item \int \sin(\ln x) dx
  \item \int_0^{\pi} e^{\cos(t)} \sin(2t) dt
  \item \int \arctan(\sqrt{x}) dx
\end{mathlist}
\begin{answer}
\begin{mathlist}
  \item 2 \sqrt{x} \sin(\sqrt{x})+ 2 \cos(\sqrt{x})+C
  \item \frac{1}{2}x \left[ \sin(\ln x)- \cos(\ln x)\right]+C
  \item \frac{4}{e}
  \item x \arctan(\sqrt{x})- \sqrt{x}+ \arctan(\sqrt{x})+C\)
      ou \((x+1)\arctan(\sqrt{x})- \sqrt{x}\)
\end{mathlist}
\end{answer}
\end{exercise}

%------------------------------------------------------------------------------%
\begin{exercise}
Decida qual método usar e calcule as integrais
(Substituição ou integração por partes).
Pode ser que você precise usar as duas técnicas no mesmo exercício.
\begin{mathlist}
\item \int \cos(x) (1+ \sin^2(x))dx
\item \int e^{x+e^x} dx
\item \int_1^3  x^4 \ln(x) dx
\item \int_0^{\frac{\pi}{6}} t \sin(2t) dt
\item \int \frac{\ln x}{x \sqrt{1+(\ln x)^2}}dx
\item \int_{-1}^1\frac{e^{\arctan(x)}}{1+x^2} dx
\item \int  \frac{x^3}{\sqrt{1+x^2}}dx
\item \int x^5 e^{-x^3} dx
\item \int_0^{1} \arcsin(x)dx
\item \int \sqrt{x}\ln(\sqrt{x})dx
\end{mathlist}
\begin{answer}
  \begin{mathlist}
  \item \sin(5)+\frac{1}{3} \sin(x)+c
  \item e^{e^x}+C
  \item \frac{243}{5} \ln 3- \frac{242}{25}
  \item -\frac{\pi}{24}+ \frac{1}{8} \sqrt{3}
  \item \sqrt{1+(\ln x)^2}+c
  \item e^{\frac{\pi}{4}}- e^{-\frac{\pi}{4}}
  \item \frac{1}{3} (1+x^2)^3-(1+x^2)+C
  \item -\frac{1}{3} e^{-x^3}(x^3+1)+C
  \item \frac{1}{2}(\pi-2)
  \item \vert x \vert \ln(\sqrt{x})- \sqrt{x}+K
  \end{mathlist}
\end{answer}
\end{exercise}

%------------------------------------------------------------------------------%
\begin{exercise}
Se $f$ é uma função contínua no intervalo $[a,b]$,
o valor médio de $f$ é definido como
\[
  f_{\text{med}}= \frac{1}{b-a} \int_a^b f(x) dx
\]
Encontre o valor médio das seguinte função
$f(x)=x \sec^2(x)$ no intervalo
$[0, \frac{\pi}{4}]$.
\begin{answer}
  $1-\frac{2}{\pi}\ln 2$
\end{answer}
\end{exercise}

%------------------------------------------------------------------------------%
%\begin{exercise}
%\begin{answer}
%\end{answer}
%\end{exercise}

%------------------------------------------------------------------------------%
%\begin{exercise}
%\begin{mathlist}[2]
%  \item
%  \item
%  \item
%\end{mathlist}
%\begin{answer}
%  \begin{alphalist}[3]
  %    \item
  %    \item
  %    \item
  %  \end{alphalist}
%\end{answer}
%\end{exercise}

\end{exercises}
%------------------------------------------------------------------------------%
